\chapter{Evaluation Protocol}
\label{ch:protocol}

All comparison maps and trajectories come from local runs on a host with two
RTX~A6000 GPUs. This chapter defines the common full-SH renderer, cameras,
alignment, and aggregation.

\begin{table}[htbp]
\centering
\small
\setlength{\tabcolsep}{4pt}
\begin{tabular}{@{}llcl@{}}
\toprule
System & Type & Rend. & Role here \\
\midrule
Ours & mono 3DGS & \checkmark & --- \\
Photo-SLAM & mono 3DGS & \checkmark & rendering \\
MonoGS & mono 3DGS & \checkmark & rendering (TUM) \\
MASt3R-SLAM & mono dense & --- & tracking control \\
VGGT-SLAM & pose+cloud & --- & tracking control \\
\bottomrule
\end{tabular}
\caption[Systems compared, and what each can produce]{Systems compared: Photo-SLAM~\cite{huang2024photoslam},
MonoGS~\cite{matsuki2024monogs}, MASt3R-SLAM~\cite{murai2025mast3rslam},
VGGT-SLAM~\cite{maggio2025vggtslam}. The latter two supply trajectory
comparisons only; MASt3R-SLAM is the inherited tracking control.}
\label{tab:systems}
\end{table}


\section{Common rendering conditions}

For each scene, $100$ frames are selected outside the union of participating
systems' keyframes. Each exported map is rendered at these same cameras by
one implementation, with identical exposure normalisation and metric code.
The resolution is $512\times288$ on Replica and $512\times384$ on TUM.
All exported spherical-harmonic bands are retained: degree three for
Photo-SLAM, degree zero for the evaluated MonoGS and our maps.

Scene PSNR is computed as $-10\log_{10}$ of the mean frame MSE, rather than
as the arithmetic mean of per-frame PSNR. LPIPS-AlexNet~\cite{zhang2018lpips}
is averaged across frames. Dataset means weight each scene equally.
Scene PSNR and the median paired frame difference can therefore rank methods
differently.

\subsection{Evaluation scope}

We call these \emph{non-keyframe reconstruction} scores. Tracked
non-keyframes may enter the refiner's supervision pool. Their exclusion from
the keyframe set therefore does not guarantee they were unseen during
optimisation. Ground-truth cameras reduce dependence on each system's own
evaluation-pose estimates, but expose tracking and alignment errors as well
as map errors. Common rendering also cannot remove differences in input
preprocessing, undistortion or optimisation history.

A strict novel-view experiment would exclude scoring views from every
optimisation and model-selection path.

\subsection{Datasets and configurations}

Rendering covers all eight Replica~\cite{straub2019replica} scenes and TUM
RGB-D~\cite{sturm2012tum} \texttt{freiburg1\_desk}, with monocular input for
all compared systems. Trajectory accuracy covers all nine TUM freiburg1
sequences. Adaptation and refinement studies additionally use
EuRoC~\cite{burri2016euroc}, ETH3D~\cite{schops2017eth3d} and
7-Scenes~\cite{shotton2013scenecoord}.

The external rendering configuration uses the released Splatt3R head,
fixed Gaussian centres, insertion thinning disabled and $300$\,s
post-sequence refinement. Photo-SLAM runs to its own shutdown schedule, so
Replica optimisation budgets are not matched. MonoGS uses a $339$\,s
colour-refinement phase on TUM desk. That phase is close in duration to our
polish, but is not a match of total pipeline compute.

Family-head adaptation and thinning use separate paired runs. Maps delivered
at sequence end are reported separately from post-sequence polish.

\subsection{Alignment}

Each monocular map has its own similarity gauge. We fit a map-to-ground-truth
alignment $(s,R,t)$ to corresponding estimated and ground-truth camera centres
using Umeyama's method~\cite{umeyama1991}. A ground-truth camera-to-world pose
is transported into map coordinates by
\begin{equation}
R_{\text{map}}=R^\top R_{\text{gt}},\qquad
t_{\text{map}}=R^\top(t_{\text{gt}}-t)/s.
\label{eq:align}
\end{equation}
The inverse scale acts on translation, not on the rotation block. Putting it
inside the rotation would produce a non-orthonormal camera orientation.
ATE is translation RMSE after similarity alignment, evaluated with
\texttt{evo\_ape tum -as}~\cite{grupp2017evo}.

MonoGS records estimated and ground-truth keyframe poses together. Its TUM
keyframes are recovered from those saved GT poses and checked against GT
orientation. This avoids confusing indices in its RGB-D-associated stream
with indices in the raw RGB sequence. The updated comparison uses this
association throughout selection, alignment and rendering.

\section{Self-consistency and provenance}
\label{sec:protocol:selfcheck}

Before interpreting a cross-system margin, the exported map is rendered at
the poses from which it was built. A poor result here requires checking
camera conventions, image association and map loading before interpreting
it as a property of the method. In an earlier MonoGS diagnostic, interpreting
camera-to-world poses as world-to-camera changed the measured score from
$18.71$ to $4.36$\,dB. Those are diagnostic values from that experiment,
not the updated common-camera scores in \cref{tab:tum}.

The current evidence records the input PLY and trajectory hashes and sizes,
alignment transforms, all scoring frame indices, per-frame metrics and
representative-view selection. It also records DC-only scores for comparison
with the older renderer, which discarded higher SH bands. The updated
Photo-SLAM Replica mean is $19.483$\,dB with full SH, compared with
$19.581$\,dB in that DC-only evaluation. The external tables use full SH
consistently; historical ablations retain their original configurations.

Qualitative views are selected at the two middle ranks of the per-frame
Ours-minus-baseline PSNR difference: MonoGS on TUM and Photo-SLAM on Replica.
All methods receive the same crop coordinates and image scale. This makes
selection repeatable without selecting the most favourable tail.

\section{Pose-source sensitivity}
\label{sec:protocol:offset}

Earlier diagnostics rendered an unchanged MonoGS map from its estimated
poses and from aligned ground-truth poses, obtaining $18.71$ and
$9.87$\,dB, respectively. A separate Photo-SLAM diagnostic measured a
smaller difference of approximately $2.6$\,dB. These experiments demonstrate
that map quality is coupled to the camera poses used to evaluate it.
They also include alignment error and the degree to which map parameters
have adapted to their training cameras.

An earlier local-to-published comparison changed views, resolution, masking,
undistortion, and optimisation budget at once. It cannot isolate a fixed
protocol offset. The comparison tables therefore contain only locally
rendered scores under the conditions above.

\section{Systems compared}
\label{sec:protocol:systems}

\Cref{tab:systems} identifies each system's available outputs.
We did not obtain a monocular MonoGS Replica map with the available
configurations; its depth-assisted configuration is excluded.

Comparable three-method TUM artefacts were obtained on desk; broader
rendering coverage remains future work. Photo-SLAM uses OpenCV~4.10 for
CUDA compatibility rather than its tested 4.7/4.8 stack.
\Cref{app:baselines} records this preprocessing-related deviation.
